\documentclass[11pt]{article}
\usepackage[a4paper,margin=1in]{geometry}
\usepackage{amsmath,amssymb,amsthm,mathtools}
\usepackage{xcolor}
\usepackage[colorlinks=true,linkcolor=blue!55!black,
            citecolor=blue!55!black,urlcolor=blue!65!black]{hyperref}

\newtheorem{theorem}{Theorem}[section]
\newtheorem{lemma}[theorem]{Lemma}
\newtheorem{remark}[theorem]{Remark}

\newcommand{\ind}{\operatorname{ind}}
\newcommand{\nul}{\operatorname{nul}}

\title{An exact reduction of the index-four conjecture\\
for the critical spherical catenoid in hyperbolic three-space}
\author{}
\date{18 August 2026}

\begin{document}
\maketitle

\begin{abstract}
\textbf{Relation to the problem list.}
This manuscript addresses Question~32 (L32), ``Every critical hyperbolic
catenoid has index four.''  Its status is \emph{partial}: it proves one global
pinching inequality and gives an exact reduction, but it does not prove the
conjecture.  The monotonicity inequality $r'(a)\ge0$ and the required
mode-two spectral inequality both remain open.

For the critical rotational free-boundary minimal annulus
$\Sigma_a\subset B^3(r(a))\subset\mathbb H^3$, $a>1/2$, Medvedev
conjectured that $\ind(\Sigma_a)=4$.  We record a short Sturm argument
showing that the geometric pinching inequality
$B(s_0(a))^2>2K(a)^2$, left as a global residual condition in recent
work, is automatic for every parameter.  We then give an exact formula
for the mode-zero index and reduce the original index-only conjecture
to two sharp scalar statements: $r'(a)\geq0$ and nonnegativity of the
first even radial eigenvalue in angular mode two.  The corresponding
strict inequalities are equivalent to the strong conclusion
$\ind=4$, $\nul=2$.  This note does not claim a global resolution:
both residual inequalities remain open, and the mode-two spectral
margin tends to zero in the large-parameter limit.
\end{abstract}

\section{Context and scope}

Medvedev proved the lower bound $\ind(\Sigma_a)\geq4$ and posed equality
as a conjecture~\cite{Medvedev}.  Pigazzini subsequently separated the
Robin Jacobi problem into Fourier modes, proved the strong conclusion
near the closing-neck limit, and reduced the remaining analysis to
one-dimensional shooting problems~\cite{PigazziniLocal,PigazziniMode}.
The purpose here is to isolate a consequence of those formulas which
removes one of the geometric hypotheses and to state exactly what is
still needed for the original, index-only conjecture.  In particular,
none of the numerical evidence mentioned below is used as a proof.

\section{An exact index reduction}

Fix $a>\frac12$ and put
\[
 B(s)^2=a\cosh(2s)-\frac12,
 \qquad K^2=a^2-\frac14,
 \qquad B_0=B(s_0(a)).
\]
The critical catenoid is parametrized on
$[-s_0,s_0]\times\mathbb S^1$, its metric is
$ds^2+B(s)^2d\theta^2$, and
$|\mathrm{II}|^2=2K^2/B^4$.  After separation into angular modes,
the radial operator in mode $k$ is
\begin{equation}\label{eq:radial-operator-reduction}
 \mathcal A_k u
 =-\frac1B(Bu')'
  +\left(2+\frac{k^2}{B^2}-\frac{2K^2}{B^4}\right)u,
\end{equation}
with $u'(s_0)=\coth(r)u(s_0)$.  On $[0,s_0]$, the even and odd
sectors have, respectively, $u'(0)=0$ and $u(0)=0$.

We use the elementary shooting count in the following form.  If $w$ is
the zero-energy solution satisfying the left boundary condition and
$w(s_0)\ne0$, then
\begin{equation}\label{eq:shooting-count-reduction}
 N_-=n_z(w)+
 \mathbf 1_{\{w'(s_0)/w(s_0)-\coth r<0\}},
\end{equation}
where $n_z(w)$ counts zeros in $(0,s_0)$.  The cases in which
$w(s_0)=0$ follow either directly from the Pruefer angle or by a
one-sided limiting argument.

\begin{lemma}[The geometric inequality is automatic]
For every $a>\frac12$,
\begin{equation}\label{eq:G-global-reduction}
 \cosh(2s_0)>2a,
 \qquad\text{equivalently}\qquad B_0^2>2K^2.
\end{equation}
The odd radial sector of mode zero has exactly one negative eigenvalue
and no zero eigenvalue.
\end{lemma}

\begin{proof}
The axial-boost Jacobi field gives the zero-energy odd solution
\[
 v(s)=\frac{a\sinh(2s)}{B(s)}>0\qquad(0<s\le s_0).
\]
It has no interior zero, and direct differentiation, together with
$\coth r=a\sinh(2s_0)/B_0^2$, gives
\begin{equation}\label{eq:odd-defect-reduction}
 \frac{v'(s_0)}{v(s_0)}-\coth r
 =\frac{2a-\cosh(2s_0)}{B_0^2\sinh(2s_0)}.
\end{equation}
On the other hand, the ambient coordinate
$\Phi^1=A(s)\sinh\varphi(s)$ is odd and satisfies
\[
 \mathcal S(\Phi^1,\Phi^1)
 =-\int_{\Sigma_a}|\mathrm{II}|^2(\Phi^1)^2\,dA<0.
\]
Thus the odd mode-zero problem has at least one negative eigenvalue.
Applying \eqref{eq:shooting-count-reduction} to $v$, whose zero count
is zero, forces the left side of \eqref{eq:odd-defect-reduction} to be
strictly negative.  This proves \eqref{eq:G-global-reduction}; the same
shooting formula then gives exactly one negative eigenvalue.  Strictness
also shows that $v$ does not satisfy the Robin condition, so the odd
kernel is trivial.
\end{proof}

Let
\[
 \phi_a=\langle\partial_a\Phi_a,\nu_a\rangle_L
\]
be the even parametric Jacobi field.  It satisfies
\begin{equation}\label{eq:parametric-data-reduction}
 \mathcal A_0\phi_a=0,
 \quad \phi_a(0)=\frac1{2K}>0,
 \quad \phi_a'(0)=0,
 \quad
 \phi_a'(s_0)-\coth r\,\phi_a(s_0)
 =-r'(a)\frac K{B_0^2}.
\end{equation}
The weighted Wronskian of $\phi_a$ and the axial-boost field is a
nonzero constant.  Sturm separation therefore shows that $\phi_a$ has
at most one zero in $(0,s_0]$.

\begin{lemma}[Exact mode-zero count]\label{lem:mode-zero-count-reduction}
The number of negative eigenvalues in the even radial sector of mode
zero is
\[
 N_-^{0,\mathrm{even}}
 =1+\mathbf 1_{\{r'(a)<0\}}.
\]
Its nullity is one if $r'(a)=0$ and zero otherwise.  Consequently the
full mode-zero contribution to the Morse index is
\begin{equation}\label{eq:mode-zero-index-reduction}
 \operatorname{ind}_0(\Sigma_a)
 =2+\mathbf 1_{\{r'(a)<0\}}.
\end{equation}
\end{lemma}

\begin{proof}
The positive timelike coordinate
$\Phi^0=A(s)\cosh\varphi(s)$ is even and
$\mathcal S(\Phi^0,\Phi^0)<0$.  Hence the even problem has at least one
negative eigenvalue.

Suppose first that $r'(a)>0$.  If $\phi_a(s_0)>0$, then $\phi_a$ has no
zero and its boundary defect divided by its boundary value is negative;
\eqref{eq:shooting-count-reduction} gives one negative eigenvalue.  If
$\phi_a(s_0)<0$, then $\phi_a$ has exactly one zero and the corresponding
boundary quotient is positive, giving again one negative eigenvalue.
If $\phi_a(s_0)=0$, then
$\phi_a'(s_0)=-r'(a)K/B_0^2<0$.  Thus the zero-energy Pruefer angle
at $s_0$ is exactly $\pi$ (the solution is positive in the open
interval and crosses zero with negative derivative).  The Robin angle
lies strictly between $0$ and $\pi$, so precisely one Robin eigenangle
has been crossed before energy zero.  Again $N_-=1$.

If $r'(a)=0$, then $\phi_a$ is a nonzero Robin eigenfunction at
eigenvalue zero.  Since the ground eigenvalue is negative and $\phi_a$
has at most one zero, zero is the second eigenvalue.  Thus there is one
negative eigenvalue and a one-dimensional kernel.

Finally suppose that $r'(a)<0$.  The alternative
$\phi_a(s_0)>0$ would give zero interior zeros and a positive boundary
quotient, hence no negative eigenvalue, contrary to the $\Phi^0$ test.
The endpoint case is also impossible: it would give
$\phi_a'(s_0)=-r'(a)K/B_0^2>0$; starting positive, the solution would
then have to be negative immediately to the left of $s_0$, producing
one interior zero in addition to the endpoint zero, contrary to the
at-most-one-zero property.  Thus $\phi_a(s_0)<0$: it has one interior
zero, while \eqref{eq:parametric-data-reduction} gives a negative
boundary quotient.  Formula \eqref{eq:shooting-count-reduction} gives
two negative eigenvalues.  Combining this even count with the preceding
odd count proves \eqref{eq:mode-zero-index-reduction} and the nullity
claim.
\end{proof}

Let $\mu_0^{\mathrm{even}}(2)$ denote the first even radial eigenvalue
of \eqref{eq:radial-operator-reduction} in mode $2$.  The mode-$1$
sector has index two and nullity two.  For every $k\ge2$, the odd radial
sector is strictly positive, all even eigenvalues except possibly the
first are positive, and
\begin{equation}\label{eq:mode-monotonicity-reduction}
 \mu_0^{\mathrm{even}}(k)
 \ge \mu_0^{\mathrm{even}}(2)+\frac{k^2-4}{B_0^2}.
\end{equation}
We have therefore proved the following exact reduction.

\begin{theorem}[Exact reduction of the index-four conjecture]
For a fixed $a>\frac12$,
\begin{equation}\label{eq:exact-index-formula-reduction}
 \operatorname{ind}(\Sigma_a)
 =4+\mathbf 1_{\{r'(a)<0\}}
  +2\sum_{k\ge2}
       \mathbf 1_{\{\mu_0^{\mathrm{even}}(k)<0\}}.
\end{equation}
In particular,
\begin{equation}\label{eq:weak-equivalence-reduction}
 \operatorname{ind}(\Sigma_a)=4
 \quad\Longleftrightarrow\quad
 r'(a)\ge0\quad\text{and}\quad
 \mu_0^{\mathrm{even}}(2)\ge0.
\end{equation}
Moreover,
\begin{equation}\label{eq:strong-equivalence-reduction}
 \operatorname{ind}(\Sigma_a)=4,
 \quad\operatorname{nul}(\Sigma_a)=2
 \quad\Longleftrightarrow\quad
 r'(a)>0\quad\text{and}\quad
 \mu_0^{\mathrm{even}}(2)>0.
\end{equation}
\end{theorem}

The remaining mode-$2$ inequality is equivalently the nonnegativity of
the quadratic form
\begin{equation}\label{eq:k2-picone-reduction}
 \mathcal S_2(Bh,Bh)
 =\int_{-s_0}^{s_0}
 \left[B^3(h')^2+\frac{3B^2-2K^2}{B}h^2\right]ds.
\end{equation}
For $a>1$ the coefficient $3B^2-2K^2$ is negative near the neck, so
\eqref{eq:k2-picone-reduction} is the genuine residual spectral
obstruction; it is not settled merely by \eqref{eq:G-global-reduction}.

\section{What remains open}

Lemma~\ref{lem:mode-zero-count-reduction} shows that the radius
monotonicity question is not merely auxiliary: a parameter with
$r'(a)<0$ would force an additional negative mode and disprove the
index-four conjecture.  Conversely, $r'(a)=0$ does not change the
index, but it creates an additional even mode-zero Jacobi field and
therefore violates the strong nullity statement.  Global monotonicity
of $a\mapsto r(a)$ is still posed as an open problem in current
work~\cite{PigazziniDegenerate}.

The second obstruction is the sign of
$\mu_0^{\mathrm{even}}(2)$.  In the conformal coordinate $x$ defined
by $dx=ds/B$, put $p=(\log B)_x=B_s$.  If $z$ is the even zero-energy
shot in mode two, normalized by $z(0)=1$ and $z_x(0)=0$, then
\[
 z_{xx}=2(1+p^2)z,
 \qquad
 \mu_0^{\mathrm{even}}(2)\geq0
 \quad\Longleftrightarrow\quad
 z_x(T)-p(T)z(T)\geq0,
\]
where $T$ is the boundary value of $x$.  Moreover,
\[
 \frac{d}{dx}\bigl(z_xB-zB_x\bigr)
 =zB\left(3-\frac{2K^2}{B^2}\right).
\]
The factor in parentheses changes sign exactly once.  Thus the
remaining assertion is a sharp cancellation inequality, not a
pointwise consequence of the global pinching lemma.

Floating-point shooting over a wide logarithmic parameter grid finds
no counterexample and finds both residual quantities positive.  It
also indicates that the mode-two boundary defect decays cubically in
$a$.  Because that margin tends to zero, a finite grid without explicit
uniform remainder bounds cannot certify the conjecture.

\begin{remark}[Status]
As of 18 August 2026, the global conjecture is still stated as open in
the primary literature.  The result proved above is a rigorous exact
reduction and a global closure of the pinching condition, not a proof
of the two remaining inequalities.
\end{remark}

\begin{thebibliography}{9}

\bibitem{Medvedev}
V.~Medvedev,
\emph{On free boundary minimal submanifolds in geodesic balls in
hyperbolic and spherical spaces},
arXiv:2311.02409v3 (2025),
\url{https://arxiv.org/abs/2311.02409}.

\bibitem{PigazziniLocal}
A.~Pigazzini,
\emph{Analytic local resolution of Medvedev's Morse index conjecture
for the critical spherical catenoid in $\mathbb H^3$},
arXiv:2605.13562v7 (2026),
\url{https://arxiv.org/abs/2605.13562}.

\bibitem{PigazziniMode}
A.~Pigazzini,
\emph{Robin nullity in mode $|k|=1$ and asymptotic radius of the
critical spherical catenoid},
arXiv:2605.11244v3 (2026),
\url{https://arxiv.org/abs/2605.11244}.

\bibitem{PigazziniDegenerate}
A.~Pigazzini,
\emph{A degenerate free boundary minimal annulus and non-uniqueness
in spherical caps beyond the hemisphere},
arXiv:2607.23534v3 (2026),
\url{https://arxiv.org/abs/2607.23534}.

\end{thebibliography}

\end{document}
